\subsection{Generic avoidance of every odd torsion order}
\label{sec:generic-all-odd}

\paragraph{Finding A21 (both positive-rank genus-two strata).}
Let $k$ be an algebraically closed field of characteristic two,
let $f\in\{1,2\}$, and let $Z$ be an irreducible component of
$\mathcal M_2^f$.  Let $C$ denote its geometric generic curve,
put $X=C^{(1)}$, and let $B$ be the locally exact differential
line bundle on $X$.  Its Raynaud theta support
\[
 \Theta=\{L\in\operatorname{Pic}^0(X):
                       H^0(X,B\otimes L)\ne0\}
\]
contains no nonzero point of odd order:
\[
 \Theta\cap
 \bigl(\operatorname{Pic}^0(X)[n]\setminus\{0\}\bigr)
             =\varnothing\qquad(n\ge1\text{ odd}).             \tag{GO.1}
\]
Consequently, for every odd $n>1$, every connected cyclic
degree-$n$ etale cover of $C$ has ordinary whole new factor.
This includes prime powers and composite degrees on the
$2$-rank-one stratum.  All assertions are unconditional.

The statement about one geometric generic curve does not
assert that a single nonempty open subset avoids all odd
orders simultaneously.  For any prescribed finite list of
odd integers the avoidance conditions do hold on a common
nonempty open, and the finite-field specialization argument
of A08 then applies for either $f=1$ or $f=2$.

\paragraph{Exact published inputs.}
Achter--Pries, \emph{Monodromy of the $p$-rank strata of the
moduli space of curves}, IMRN 2008, Art.~rnn053,
Theorem~4.5 (author version p.~14), gives full
$\operatorname{Sp}_{2g}(\mathbf F_\ell)$ and
$\operatorname{Sp}_{2g}(\mathbf Z_\ell)$ monodromy for
$\ell\ne p$, including $g=2$ and $f=1,2$
\cite{AchterPries2008}.
For the generic-point interpretation used here, the proof
of Ozman--Pries, Proposition~4.3, printed p.~465, explicitly
states that the full symplectic level-$\ell$ torsor $Z[\ell]$
is irreducible \cite[Proposition~4.3, proof]{OzmanPries2019}.
Thus this input is stronger than mere
connectedness of a cover over a possibly nonnormal base.
The tricanonical framing construction in Achter--Pries,
Section~5.1 and Lemma~5.2, pp.~17--18, passes to a scheme
without reducing monodromy.  The lifting assertion for a
closed subgroup of $\operatorname{Sp}_4(\mathbf Z_\ell)$
surjecting modulo $\ell$ is also
Landesman--Swaminathan--Tao--Xu,
\emph{Lifting subgroups of symplectic groups over
$\mathbf Z/\ell\mathbf Z$}, Research in Number Theory \textbf{3}
(2017), Art.~14, Theorem~1; their Proposition~1(b) says that
$\operatorname{Sp}_4(\mathbf Z/\ell^a)$ is perfect for every
odd $\ell$ and $a\ge1$ \cite{LSTX2017}.  Finally, Hall,
\emph{Classical Groups and Geometry}, version 6 May 2015,
Theorem~12.26, printed p.~173, gives the center and simplicity
of $\operatorname{PSp}_4(\mathbf F_\ell)$ for all odd $\ell$,
including $\ell=3$ \cite{Hall2015}.

\begin{proof}
\emph{1. The monodromy statement at the geometric generic
point.}
Use tricanonical frames to replace $Z$ by a scheme carrying
its universal curve.  The frame morphism has geometrically
integral $\operatorname{GL}_5$ fibres.  In particular its
pullback to the irreducible full level-$\ell$ torsor
$Z[\ell]$ is irreducible.  Restrict to a nonempty smooth
affine open $T$ of the resulting integral scheme.  The
level-$\ell$ torsor remains irreducible on $T$, so its
geometric generic Galois image is the full
$\operatorname{Sp}_4(\mathbf F_\ell)$.
The closed $\ell$-adic image over this same generic field
is consequently $\operatorname{Sp}_4(\mathbf Z_\ell)$ by
the lifting theorem just cited.  This reasoning works for
every odd $\ell$ on the same generic curve.  It does not
assume that deleting a singular locus from an arbitrary
irreducible base preserves its fundamental group.

These statements apply equally to $J=\operatorname{Pic}^0(X)$.
Indeed relative Frobenius induces an isogeny between the
Jacobians of $C$ and $X$ of degree a power of two.  It is an
isomorphism on every odd-order Tate module, and is a
similitude of the corresponding alternating pairings.
Conjugating a full symplectic group by a similitude leaves
that symplectic group unchanged.

\emph{2. Independence at distinct odd primes.}
Write $n=\prod_{\ell\mid n}\ell^{a_\ell}$ and
$G_\ell=\operatorname{Sp}_4(\mathbf Z/\ell^{a_\ell})$.
Each $G_\ell$ is perfect.  Its reduction kernel modulo
$\ell$ is an $\ell$-group, as is seen by filtering the
congruence kernel by successive powers of $\ell$.
If $G_\ell$ has a nontrivial simple quotient, that quotient
is nonabelian, since $G_\ell$ is perfect.  The image of the
congruence kernel is a normal $\ell$-subgroup of this simple
quotient and must be trivial.  The quotient therefore
factors through $\operatorname{Sp}_4(\mathbf F_\ell)$.
Its center has trivial image in a nonabelian simple group,
and Hall's theorem now implies that the simple quotient is
exactly $\operatorname{PSp}_4(\mathbf F_\ell)$.
The orders
\[
 \bigl|\operatorname{PSp}_4(\mathbf F_\ell)\bigr|
       =\frac{\ell^4(\ell^2-1)(\ell^4-1)}2                  \tag{GO.2}
\]
strictly increase with the positive integer $\ell$.
Thus factors belonging to distinct odd primes have no
common nontrivial simple quotient.

For clarity, this implies the required product assertion
without an independence assumption.  If a subgroup
$H\subset A\times B$ surjects to $A$ and $B$, Goursat's
argument identifies it with the fibre product over
isomorphic quotients $A/N_A\simeq B/N_B$.  If $H$ is
proper, this common quotient is nontrivial and has a
nontrivial simple quotient.  This would be a common simple
quotient of $A$ and $B$.  Induct on the number of prime
factors.  A nonabelian simple quotient of a product of
perfect groups factors through one factor: the images of
the factors are commuting normal subgroups, and at most
one can be the whole nonabelian simple group.  There are
no abelian simple quotients.  The obstruction just
described is therefore absent at each induction step.

The generic Galois image on $J[n]$ surjects to every
$G_\ell$ by Step~1 and is contained in their product.
Consequently it is
\[
            \prod_{\ell\mid n}G_\ell
                  =\operatorname{Sp}_4(\mathbf Z/n).         \tag{GO.3}
\]
The equality on the right uses the Chinese remainder
decomposition of matrices and their defining symplectic
equations.

\emph{3. Transitivity and the genus-two intersection bound.}
Over the local ring $R=\mathbf Z/\ell^a$, a vector of exact
order $\ell^a$ is primitive.  If $v$ is primitive, the
perfect alternating form admits $w$ with
$\langle v,w\rangle=1$.  The span of $v,w$ is a
nondegenerate free direct summand; its orthogonal
complement is a free rank-two symplectic module.
Completing $v,w$ to a symplectic basis proves that
$\operatorname{Sp}_4(R)$ is transitive on primitive
vectors.  Applying this in each Chinese remainder factor
proves that the group in (GO.3) is transitive on the
points of exact order $n$.

Because $B$ has degree one, $\Theta$ is the translated
Abel curve
\[
                 R\longmapsto\mathcal O_X(R)\otimes B^{-1}.
                                                               \tag{GO.4}
\]
For every nonzero $T\in J$, it satisfies
\[
                   \#\bigl(\Theta\cap(\Theta+T)\bigr)\le2.
                                                               \tag{GO.5}
\]
This is the bound proved in A19.  A short proof is included
to specify its exact scope.  Three distinct intersection
points give pairs $(R_i,S_i)$ with $[R_i-S_i]=T$.
The $R_i$ and the $S_i$ are separately distinct.  For
$i=2,3$, the distinct effective divisors
$R_1+S_i\sim R_i+S_1$ move in a degree-two linear system.
Riemann--Roch in genus two makes this class canonical.
Both $S_2$ and $S_3$ would then be the unique effective
degree-one divisor in the class $K_X-R_1$, a contradiction.
If the two displayed degree-two divisors were equal as
divisors, their equality and $R_1\ne R_i$ would instead
force $T=0$.  This also is excluded.

Suppose now that $\Theta$ contains a point of exact odd
order $n>1$.  The line bundle $B$, and hence its theta
support, is defined over the generic field of the family.
Galois invariance and (GO.3) force $\Theta$ to contain
every point of exact order $n$.  Choose coordinates
$J[n]\simeq(\mathbf Z/n)^4$, let $e_1,e_2$ be two
distinct coordinate vectors, and choose an odd prime
$\ell\mid n$.  Set
\[
                   P=e_1,\qquad T=(n/\ell)e_2.
\]
The point $T$ has exact order $\ell$ and all
$P+jT$, $0\le j<\ell$, have exact order $n$, since their
first coordinate is $1$.  They therefore lie on $\Theta$.
Each also lies on $\Theta+T$, using the preceding point
in this cyclic list.  There are $\ell\ge3$ such distinct
points, contradicting (GO.5).  This excludes every exact
odd order and proves (GO.1).

\emph{4. The whole new factor of an odd cyclic cover.}
Work over the algebraic closure of the generic field.
Let $\pi:D\to C$ be a connected cyclic etale cover of
odd degree $n$.  Choose a line bundle $M$ of exact order
$n$ defining the cover, so that
\[
 \pi_*\mathcal O_D=\bigoplus_{i=0}^{n-1}M^i,\qquad
 H^1(D,\mathcal O_D)=\bigoplus_{i=0}^{n-1}H^1(C,M^i).
                                                               \tag{GO.6}
\]
For $i\ne0$, the latter summand has dimension one.
Put $L_i=(M^i)^{(1)}$ on $X=C^{(1)}$.  It is nontrivial
of odd order and $F_{C}^*L_i=(M^i)^2$.  Tensoring
$0\to\mathcal O_X\to F_{C*}\mathcal O_C\to B\to0$
by $L_i$ and using the projection formula gives
\[
 0\longrightarrow L_i\longrightarrow F_{C*}(M^i)^2
                 \longrightarrow B\otimes L_i
                 \longrightarrow0.                          \tag{GO.7}
\]
By (GO.1), $H^0(X,B\otimes L_i)=0$.  Thus
$H^1(X,L_i)\to H^1(C,M^{2i})$ is injective, hence
bijective between its one-dimensional source and target.
Under Frobenius-semilinear identification this is precisely
the arrow from character $i$ to character $2i$ in (GO.6).
Multiplication by two permutes all nonzero residues
modulo $n$, so Frobenius is bijective on the entire
nontrivial-character sum.

Its stable rank is therefore $n-1$; the trivial-character
sum has stable rank $f$.  Consequently
$f(D)=f+n-1$.  The norm--pullback isogeny decomposition
$J_D\sim J_C\times\operatorname{Prym}(D/C)$ and
Riemann--Hurwitz give
\[
 f\bigl(\operatorname{Prym}(D/C)\bigr)=n-1
       =\dim\operatorname{Prym}(D/C),
\]
which proves the final assertion.
\end{proof}

\paragraph{Scope relative to the earlier results.}
For prime $n$, the conclusion agrees with Ozman--Pries,
Theorem~1.1; for the ordinary stratum, Yang's
new-ordinarity theorem already supplies the arbitrary
odd-degree conclusion used in A08.
The argument here is a direct genus-two consequence of
published full monodromy and the intersection bound.  It
does not require an extension of the compact-type
degeneration argument from prime to composite degrees.
The full joint monodromy in (GO.3) also supplies the finite
level needed in A22 below.
